\documentclass[11pt,a4paper]{article}

\usepackage[
  a4paper,
  margin=1cm
]{geometry}

\usepackage[utf8]{inputenc}
\usepackage[T1]{fontenc}
\usepackage{lmodern}
\usepackage{hyperref}

\setlength{\parindent}{0pt}
\setlength{\parskip}{8pt}

\hypersetup{
  colorlinks=true,
  linkcolor=black,
  urlcolor=blue
}

\begin{document}

\begin{center}
  {\Large \textbf{Base Models and the AI Harness}}\\[4pt]
  {\large \textbf{A Simple Guide to How AI Models Work and What We Are Building}}
\end{center}

\vspace{0.2cm}
\hrule
\vspace{0.3cm}

\section*{1. The Base Model: A Passive Brain}

A base AI model (like ChatGPT, DeepSeek, or Qwen) is not an independent robot or an active program. At its core, it is simply a text predictor. It looks at the words you give it and guesses what the next word should be, one word at a time. Once it finishes writing its answer, it stops completely and waits for you to talk to it again.

Because of this design, a base model is completely passive. It cannot wake up on its own, it does not know what time it is, and it cannot touch the real world. It cannot open a folder on your computer, it cannot read a file, and it cannot click a button. Also, because it only guesses words based on patterns, it is very bad at doing precise math in its head and often makes up wrong numbers.

A good way to picture a base model is a powerful car engine sitting on a garage floor. The engine has a lot of horsepower, but because it has no car frame, wheels, steering wheel, or gas pedal, it cannot move anywhere by itself.

\section*{2. Why We Need an AI Harness}

Think about how a computer works. The processor (CPU) does the raw computing, but you cannot use a CPU without an Operating System like Windows or Linux. The Operating System manages memory, opens files, and lets you use a mouse and keyboard. In the same way, a base AI model needs an AI harness to do useful work.

Real-world jobs cannot be finished with a single quick answer. For example, if you ask an AI to read a scanned report, calculate remaining equipment life, and create an approval document, a raw model cannot do that alone. It has no hands to open the file and no calculator to verify the math.

An AI harness is the software that surrounds the base model and gives it hands, memory, and tools. The harness breaks a big job down into a list of smaller steps. When the AI needs information, the harness reads the file and feeds it to the model. When the AI needs to do math, the harness makes it write real code, runs that code safely, and brings back the exact numerical answer. The harness also makes sure the AI stays safe and does not delete or change important computer files.

\section*{3. What We Are Making}

We are building an AI Agent Harness that runs completely on a local computer without needing the internet.

Our harness takes a user's request and plans out the work step by step. Instead of forcing one model to do everything, our system automatically picks the best model for each task. It sends coding work to a model trained for code, pictures and scanned drawings to a model that can see images, and deep reasoning problems to a model trained for logic.

To stop the AI from making math mistakes, our harness uses a safe, locked code runner. When an engineering formula or calculation is needed, the AI writes a short Python script. Our harness runs this script inside an isolated box, gets the exact mathematical result, and passes it back to the AI.

Finally, our harness takes the AI's findings and turns them into real files that people can use right away. It can automatically write a formatted Word document with proper headings and tables, or create an Excel spreadsheet with working formulas.

In simple terms: the base model does the thinking, but our harness provides the eyes, the hands, the memory, and the tools to turn that thinking into finished work.

\end{document}
